\documentclass[aspectratio=169]{beamer}

\usetheme{Madrid}

\definecolor{NavyTIPE}{RGB}{0,42,92}
\definecolor{CyanTIPE}{RGB}{0,160,200}
\definecolor{LightCyan}{RGB}{220,244,252}

\setbeamercolor{palette primary}   {bg=NavyTIPE, fg=white}
\setbeamercolor{palette secondary} {bg=CyanTIPE, fg=white}
\setbeamercolor{palette tertiary}  {bg=NavyTIPE, fg=white}
\setbeamercolor{palette quaternary}{bg=NavyTIPE, fg=white}
\setbeamercolor{structure}{fg=NavyTIPE}
\setbeamercolor{title}{fg=white, bg=NavyTIPE}
\setbeamercolor{frametitle}{fg=white, bg=NavyTIPE}
\setbeamercolor{block title}{bg=CyanTIPE, fg=white}
\setbeamercolor{block body}{bg=LightCyan, fg=black}
\setbeamercolor{alerted text}{fg=CyanTIPE}

\setbeamertemplate{bibliography item}{}

\usepackage[utf8]{inputenc}
\usepackage[T1]{fontenc}
\usepackage[french]{babel}
\usepackage{csquotes}

\usepackage{amsmath, amssymb}
\usepackage{graphicx}
\usepackage{hyperref}
\usepackage{array}
\usepackage{colortbl}

\usepackage{tikz}
\usetikzlibrary{arrows.meta, decorations.pathmorphing, decorations.markings,
decorations.pathreplacing, patterns, calc, positioning, shapes.geometric}

\title[TIPE]{%
  Influence de la rigidité d'une hélice de drone\\
sur l'efficacité propulsive}
\subtitle{Oral d'avancement}
\author{Merlin Claret}
\institute{PTSI — Lycée la Martinière Monplaisir, Lyon}
\date{20 mars 2026}

\begin{document}

\begin{frame}[plain]
  \titlepage
\end{frame}

\begin{frame}{Plan}
  \tableofcontents
\end{frame}

\section{Contexte et lien au thème}

\begin{frame}{Contexte et enjeux}
  \begin{itemize}
    \item Les drones ont de nombreuses applications : applications
      civiles, industrielles et de loisir.
    \item L’impression 3D permet de fabriquer rapidement des pièces
      sur mesure. Les \alert{paramètres d'impression affectent la rigidité} (Palmer \& Laliberté, 2023).
    \item Cependant, les \alert{hélices FDM (Impression 3D) se déforment significativement}
      (Zhu et al., 2021) sous charge aérodynamique.
    \item Cette déformation nuit-elle réellement
      à l'efficacité propulsive~?
  \end{itemize}
\end{frame}

\begin{frame}{Lien avec « Optimisation, Sobriété, Efficacité »}
  \begin{description}[\textbf{Optimisation}]
    \item[\textcolor{NavyTIPE}{Sobriété}]
      Limiter la masse de matière (épaisseur, remplissage, choix du
      polymère).
    \item[\textcolor{NavyTIPE}{Efficacité}]
      Maximiser le rapport poussée\,/\,puissance en tenant
      compte de la \alert{déformation aéroélastique}.
    \item[\textcolor{NavyTIPE}{Optimisation}]
      Recherche d'un \alert{compromis de rigidité} — ni trop souple
      (perte aérodynamique), ni trop rigide (surcoût matière).
  \end{description}
\end{frame}

\section{Problématique}

\begin{frame}{Formulation de la problématique}
  \begin{block}{Problématique}
    \textit{Quelle est l'influence du module d'Young effectif d'une hélice de
      drone imprimée en 3D sur sa déformation en rotation et sur la perte
    d'efficacité propulsive~?}
  \end{block}

  \vspace{0.2cm}

  \begin{itemize}
    \item Comparaison de plusieurs matériaux (PLA, PETG, composites)
      et de methodes d'impression.
    \item
      \alert{rigidité mesurée} $\;\rightarrow\;$
      \alert{déformation en rotation} $\;\rightarrow\;$
      \alert{variation de poussée}.
  \end{itemize}
\end{frame}

\section{Pistes de modélisation}

\begin{frame}{Modélisation du comportement de l'hélice}
  \textbf{\textcolor{NavyTIPE}{Aspect mécanique (Flexion)}}
  \begin{itemize}
    \item Assimilation de la pale à une poutre encastrée.
    \item \alert{Déformation mécanique de la pale} conditionnée par la rigidité du matériau (Zhu et al., 2021).
    \item Plus l'hélice est souple, plus elle fléchit sous la charge de l'air.
  \end{itemize}

  \vspace{0.3cm}
  \textbf{\textcolor{NavyTIPE}{Aspect aérodynamique (Portance)}}
  \begin{itemize}
    \item La flexion de la pale modifie son angle d'attaque.
    \item Conséquence directe : la portance et la traînée locales changent.
    \item But : mesurer l'écart de performance.
  \end{itemize}

  \vspace{0.4cm}
  \begin{block}{Stratégie de modélisation}
    Observer l'interaction entre les forces de l'air et les déformations. \\
    \textit{On peut utiliser le modèle de poutre encastrée d'Euler--Bernoulli pour prédire la déformation, et le valider expérimentalement.}
  \end{block}
\end{frame}

\section{Piste démarche expérimentale}

\begin{frame}{Mesure de la rigidité des hélices}
  \begin{itemize}
    \item Impression d'éprouvettes avec des \alert{paramètres de remplissage identiques} à ceux des hélices.
    \item \alert{Mesure du module d'Young en traction} pour déterminer la rigidité de l'impression.
  \end{itemize}
  \vspace{0.4cm}

  \begin{block}{But de la manipulation}
    Ne pas se fier uniquement au données du matériau, mais obtenir la rigidité de la configuration réelle d'impression pour l'associer ensuite à la déformation de l'hélice en vol.
  \end{block}
\end{frame}

\begin{frame}{Banc de poussée et déformation}
  \begin{columns}
    \column{0.5\textwidth}
    \begin{itemize}
      \item Montage : moteur + hélice sur un \alert{banc de poussée}.
      \item Instrumentation prévue :
        \begin{itemize}
          \item Balance (mesure de la poussée)
          \item Tachymètre (mesure du régime RPM)
          \item Caméra rapide (suivi de la déformation de l'hélice)
        \end{itemize}
    \end{itemize}

    \column{0.5\textwidth}
    \centering
    \begin{tikzpicture}[scale=0.8, every node/.style={transform shape}]
      % Table / Support
      \draw[thick] (-0.5, 0) -- (4.5, 0);
      \fill[pattern=north east lines] (-0.5, -0.3) rectangle (4.5, 0);

      % Balance
      \draw[thick, fill=CyanTIPE!20] (1.5, 0) rectangle (2.5, 0.5) node[midway] {\footnotesize Balance};

      % Support moteur
      \draw[thick, fill=gray!20] (1.8, 0.5) rectangle (2.2, 2.5);

      % Moteur
      \draw[thick, fill=NavyTIPE!80] (1.6, 2.5) rectangle (2.4, 3) node[midway, white] {\footnotesize Moteur};

      % Hélice
      \draw[thick, fill=CyanTIPE] (0.5, 3) ellipse (1.5cm and 0.1cm);
      \node[above] at (0.5, 3.1) {\footnotesize Hélice};

      % Flux d'air
      \draw[->, thick, CyanTIPE, dashed] (0.8, 3.5) -- (0.8, 4.2);
      \draw[->, thick, CyanTIPE, dashed] (1.5, 3.5) -- (1.5, 4.2);
      \node[right, CyanTIPE] at (1.5, 3.8) {\footnotesize Poussée};

      % Caméra
      \node[draw, thick, fill=gray!30, minimum width=0.8cm, minimum height=0.5cm] (cam) at (-1, 2.5) {};
      \draw[thick] (-0.6, 2.5) -- (-0.3, 2.7) -- (-0.3, 2.3) -- cycle;
      \node[below] at (-1, 2.2) {\footnotesize Caméra};
      \draw[dashed, gray] (-0.3, 2.7) -- (1.5, 3.2);
      \draw[dashed, gray] (-0.3, 2.3) -- (1.5, 2.8);
    \end{tikzpicture}
  \end{columns}
\end{frame}

\begin{frame}{Suivi visuel de l'hélice}
  \begin{itemize}
    \item Observation de la pale de \alert{profil} pendant la rotation.
    \item Utilisation de marques sur le bout de l'hélice pour repérer sa position.
    \item Pointage depuis la vidéo pour \alert{estimer la flexion} subie par la pale.
    \item Obtenir une preuve directe que l'hélice se plie, et mettre ce pliage en relation avec la poussée.
  \end{itemize}
\end{frame}

\section{Références}

\begin{frame}[allowframebreaks]{Références}
  \small
  \begin{thebibliography}{99}

    \bibitem{Palmer2025}
    Palmer, M., \& Laliberté, J. (2023). Effects of non-planar slicing techniques and carbon fibre material additives on the mechanical properties of 3D-printed drone propellers. \textit{Drone Systems and Applications}, \textit{11}.

    \bibitem{AerodynamicUAV}
    Zhu, H., Jiang, Z., Zhao, H., Pei, S., Li, H., \& Lan, Y. (2021). Aerodynamic performance of propellers for multirotor unmanned aerial vehicles: Measurement, analysis, and experiment. \textit{Shock and Vibration}, \textit{2021}.

    \bibitem{Forest2022}
    Forest, S., \& Amestoy, M. (2022). \textit{Mécanique des milieux continus. Volume 2, Pratique}. Presses des Mines.

  \end{thebibliography}
\end{frame}

\end{document}
